% SPDX-License-Identifier: CC-BY-4.0
% Copyright (c) 2026 Martin Schröder <info@swedishembedded.com>

\begin{abstract}
A local model that is meant to keep learning from the documents, tools and repositories
it works with needs a loop from source to release that can be inspected stage by stage.
We describe the self-learning flywheel built by three cooperating projects: Splinter decides
what is learned, Brain trains and serves, and Sven runs the agents. The loop takes a source,
has a model propose tasks that code admits or refuses for named reasons, attempts them
through an agent runtime, keeps what typed verifiers confirm, projects verified experience
into training data, trains a LoRA adapter or consolidates into the weights, and releases the
result only if four checks over unseen material pass, with immutable releases and rollback.
Every statement about the system carries an id in a claims ledger and is re-checked against
the three repositories at pinned commits (116 structure claims, 223 checks); every number
comes from a raw run record. Four small experiments exercise the loop. A task funnel on three
documentation pages admitted 63 tasks and refused 84 proposals, 82 of them because the
generator did not name the task's subject. A two-generation turn on invented facts stopped before
the release gate in both generations, because training refused a split of 8 and of 4
records, so no improvement, retention or rollback is measured; the unadapted policy answered
none of the invented facts. Qualification runs of the 27B model on one 95~GiB card fit an
NF4 base at 8192 tokens in about 36~GiB and a bf16 base at 16384 tokens in about 78~GiB.
We list what was not measured.
\end{abstract}
